\SetPractica{Curva de calibración de nitritos}{Curva de calibración de nitritos}

\section{Introducción teórica}



\section{Objetivos}

\begin{itemize}
    \item 
\end{itemize}

\section{Materiales, equipos y reactivos}

\begin{table}[H]
\centering{\small\setstretch{1.25}
\begin{tabular}{|m{0.3\linewidth}|m{0.3\linewidth}|m{0.3\linewidth}|} \hline
    
    \thead{Equipos} & \thead{Materiales} & \thead{Reactivos} \\ \hline
    
    {Espectrofotómetro \par} 
    
    &
    
    {Vasos de precipitado \par}
    {Micropipeta 1mL \par}
    {Micropipeta 2mL \par}
    {Puntillas para micropipetas \par}
    {Matraz volumétrico \par}
    {Embudo \par}
    {Pipetas graduadas \par}
    {Pipetas volumétricas \par}
    {Celdas para espectrofotómetro}
    
    & 
    
    {Agua destilada \par}
    {Reactivo color-buffer \par}
    {Disolución de referencia de trabajo de nitritos \par}
    {Tiras de pH}
    
    \\ \hline
\end{tabular}}\end{table}

\section{Procedimientos}

\begin{enumerate}
    \item En matraces volumétricos de 50mL preparar las disoluciones de materiales de referencia (\autoref{tab:disoluciones_de_materiales_de_referencia})
    \item Añadir 2mL de reactivo de color-buffer para cada disolución de referencia; mezclar y permitir que el color se desarrolle en por lo menos 15 min. El medio de la reacción para la coloración debe tener un pH entre 1,5 y 2,0 unidades de pH, esta medición puede realizarse con papel indicador. 
    \item Leer la absorbancia en el espectrofotómetro a 540 nm ó 543 nm.
    \item Con ayuda de una hoja de cálculo, elaborar los cálculos pertinentes para trazar la curva de   calibración y comprobar los controles de calidad r $\geq$ 0.998 y un porcentaje de recobro entre 80-120\%
\end{enumerate}

\begin{table}[H]
\caption{Disoluciones de materiales de referencia}
\label{tab:disoluciones_de_materiales_de_referencia}
\centering{\small\setstretch{1.25}
\begin{tabular}{|m{0.485\linewidth}|m{0.485\linewidth}|}
        \thead{Volumen en mL de disolución de referencia de trabajo de nitritos con concentración de 1mg/L de nitritos}
        &
        \thead{Concentración de masa final de en mg/L de nitritos cuando se diluye a 50mL}
        \\ \hline
        
        0   & Blanco \\
        1   & 0.02   \\
        1.5 & 0.03   \\
        2   & 0.04   \\
        5   & 0.1    \\
        10  & 0.2    \\
        
    \hline        
\end{tabular}}\end{table}